\section{Preliminaries}\label{sec:aggregation}
\paragraph{Structures and formulas.}
A vocabulary is a finite set of binary relation symbols. A structure $\A$
interprets each symbol $R$ as a binary relation $R^{\A}$ on a nonempty
domain $D$; all structures in this paper are finite. \FOthree\ is
first-order logic with equality restricted to the three variable names
$x,y,z$; a name may be quantified repeatedly. For a formula $\varphi$ whose
free variables are among $x,y$, write $\varphi^{\A}\subseteq D^2$ for the
relation that it defines in $\A$ and $\varphi^{\A}(u,v)\in\{0,1\}$ for its
value
at a pair. A formula with at most two free variables can be brought to this
form by permuting the three names, which changes neither its size nor the
relation it defines. The size $|\varphi|$ is the number of nodes of the
syntax tree,
so that $|\Psi_k|=\Theta(k)$ and $|\Phi_k|=\Theta(k)$.

\paragraph{The calculus of relations.}
\CoR\ terms are built from the relation symbols and the constants
$\zero=\varnothing$, $\one=D^2$, and $\id=\Delta_D$ (the identity relation) by
union, intersection, complement relative to $D^2$, converse
$P^\smile=\{(v,u):(u,v)\in P\}$, and composition
\[
 (P;Q)(u,v)\iff\exists z\in D\,[P(u,z)\land Q(z,v)].
\]
The endpoints and the witness need not be distinct. We write $t^{\A}$ for
the relation defined by $t$ in $\A$, $|t|$ for the number of nodes of the
syntax tree of $t$, and $\cost{t}$ for its number of composition nodes. In a
syntax tree all of whose nodes have arity at most two, the leaves outnumber
the binary nodes by exactly one; since the compositions are among the
binary nodes, a term has at least $\cost{t}+1$ leaves, and counting the
leaves and the compositions alone already gives
\begin{equation}\label{eq:tree-size}
 |t|\ge2\cost{t}+1.
\end{equation}
A \emph{relational circuit} $\Comp$ is a directed acyclic graph whose gates
carry the same operations, whose sources are relation symbols and constants,
and which has one designated output gate. It is evaluated in topological
order, and an intermediate relation may be used by several gates. We write
$\gcost{\Comp}$ for the number of composition gates and $|\Comp|$ for the
\emph{size}, the number of nodes of the graph, sources and gates together;
this is the analogue for circuits of the node count $|t|$ of a term.
Neither the depth nor the number of gates of other kinds is bounded, so the
only relation between the two measures is $|\Comp|\ge\gcost{\Comp}$. A term
is a circuit
in which every gate is used once, so every statement about circuits below
applies to terms.

A term or circuit is fixed once and for all and is evaluated on every finite
structure over its vocabulary. It contains no constants beyond $\zero$,
$\one$ and $\id$, which are interpreted in each structure; in particular, it
contains no constant chosen for a particular domain or domain size.
Section~\ref{sec:robustness} shows that the lower bounds persist even when
the circuit may depend on the domain size. We write
$t\equiv_{\rm fin}\varphi$ if $t^{\A}=\varphi^{\A}$ for every finite
structure $\A$, and use the same notation for circuits.

\paragraph{Notation.}
$[k]=\{1,\ldots,k\}$, $2^{[k]}$ is the set of all subsets of $[k]$ and
$\binom{[k]}{r}$ the set of its $r$-element subsets. Blackboard bold is
reserved for number systems and for scalars: $\mathbb N$, $\mathbb Z$,
$\mathbb R$, $\mathbb C$, the field $\mathbb F_q$, the truth values
$\bool$, and $\mathbb S$ for a semiring; these never clash with the
italic $N$ for a number of colors or with $\Comp$ for a computation. From
Section~\ref{sec:preservation} on, $u,v$ denote vertices of the finite
structure, $z$ a witness, $h$ its hub and $x$ one of its anchors. Colors are
written $\lambda,\mu,\nu$ where the set $\Col$ is left abstract, and $A$ or
$B$ where $\Col$ is a concrete family of subsets of $[k]$;
Section~\ref{sec:preservation} uses only the first convention and
Section~\ref{sec:logical} only the second, while
Sections~\ref{sec:robustness} and~\ref{sec:matrix} use both. Cells are
written
$C_\lambda$, $C_J$, $C_\Col$,
always with a subscript, so that the unsubscripted $\Comp$ is unambiguously
a circuit or a computation. The letters $P,Q$ without a subscript are the
two operands of an aggregation; $P_i,Q_i$ with a subscript are the input
symbols of \eqref{eq:phi}. For
$\varnothing\ne D'\subseteq D$, the substructure $\A\res D'$ has domain $D'$
and relations $R^{\A}\cap D'^2$, and for a relation $P$ on $D$ we write
$P\res D'=P\cap D'^2$. Every \CoR\ operation commutes with this restriction
except composition, whose witnesses are confined to $D'$:
$(P;Q)\res D'\supseteq(P\res D');(Q\res D')$, and the inclusion can be
strict. Controlling this defect is the subject of the next section.
